\documentclass[11pt]{article}
\usepackage[margin=1in]{geometry}
\usepackage{amsmath,amssymb,amsthm,mathtools,booktabs,array,microtype}
\usepackage[hidelinks]{hyperref}
\newtheorem{theorem}{Theorem}
\newtheorem{lemma}[theorem]{Lemma}
\newtheorem{proposition}[theorem]{Proposition}
\newcommand{\C}{\mathbb C}
\newcommand{\R}{\mathbb R}
\newcommand{\Q}{\mathbb Q}
\newcommand{\Z}{\mathbb Z}
\newcommand{\tr}{\operatorname{tr}}
\newcommand{\rank}{\operatorname{rank}}
\newcommand{\diag}{\operatorname{diag}}
\newcommand{\one}{\mathbf 1}
\newcommand{\pathcode}[1]{\nolinkurl{#1}}
\title{Projective asymmetry and exact moment feasibility for nine equiangular tight lines in $\C^4$}
\author{}
\date{}
\begin{document}
\maketitle

\begin{abstract}
Any complex equiangular tight frame of nine lines in dimension four has trivial projective unitary automorphism group and admits no antiunitary symmetry. We exclude the nine prime-order unitary permutation types by trace and cyclic-span arguments together with a Fourier obstruction for three free orbits. For antiunitaries, a real Gram reduction leaves three cases. One admits an explicit rational sum-of-squares identity, and the other two admit factorized polynomial certificates with rigorously bounded residuals. In the unrestricted problem, we give exact feasible moments for the full quartic relaxation and for rank-strengthened sextic relaxations with three through six vertices in the cubic monomial supports. The five-vertex problem has zero common spectral margin before a second facial reduction and admits strictly feasible moments on the reduced face. The six-vertex relaxation tests all cubic monomials with a selected family of degree-six equalities and has a strict witness after structural facial reduction. Existence of the frame remains open.
\end{abstract}

\section{Introduction}

An equiangular tight frame (ETF) is a spanning family of unit vectors whose pairwise inner products have constant modulus and whose frame operator is a scalar multiple of the identity. It attains equality in the Welch coherence bound \cite{welch}. Write $\operatorname{ETF}(d,N)$ for $N$ distinct such lines in $\C^d$. Naimark complementation exchanges $d$ and $N-d$. We consider representatives with $N>2d$, using the smaller dimension in each such pair. The complex absolute bound $N\le d^2$ rules out $d\le2$ in this range. For $d=3$ the only possibilities are $N=7,8,9$. Constructions exist for $(3,7)$ and $(3,9)$, and Sz\"oll\H{o}si excludes $(3,8)$ and its Naimark complement $(5,8)$ by a Gr\"obner-basis computation. Thus $(4,9)$ is the smallest unresolved pair by $N$ in the regime $N>2d$, up to Naimark complementation, in the literature we checked \cite[pp.~3, 5 and Table~2]{fickusmixon}\cite[Proposition~4.2 and Theorem~4.3]{szollosi}.

The $\operatorname{ETF}(4,9)$ parameters are absent from the existence table of Fickus and Mixon \cite[Table~2]{fickusmixon}. The numerical packings of Dhillon, Heath, Strohmer, and Tropp and the spectral searches of Goyeneche and Turek give benchmarks for small parameters \cite[Table~2]{dhillon}\cite[Sections~VI--VII]{goyeneche}. The author-maintained packing table of Jasper, King, and Mixon lists coherence $0.40185012$ for its best submitted nine-line packing in dimension four, above the Welch value $\sqrt{10}/8\approx0.39528471$. A submitted packing supplies an upper bound on optimal coherence \cite{gameofsloanes}. The Singer--Zauner gap in dimension four excludes $N=14,15$ \cite{gap}.

Projective unitary symmetry acts on the lines underlying a frame, so it allows each unit representative to acquire a phase. This is the symmetry notion associated with projective unitary equivalence and the corresponding cyclic Gram products \cite{chienwaldron}. We also consider conjugate-linear isometries. The two main results constrain the symmetries of a putative frame and establish feasibility of specific relaxations of the unrestricted Gram equations.

\begin{theorem}[Projective asymmetry]\label{thm:asymmetry}
Suppose nine lines in $\C^4$ admit unit representatives forming an ETF. Every unitary preserving the line set is a scalar, and no antiunitary preserves the line set.
\end{theorem}

\begin{theorem}[Exact moment feasibility]\label{thm:moment}
The anchored Gram equations for a complex $\operatorname{ETF}(4,9)$ have feasible moments in the full quartic relaxation and in the rank-strengthened three-, four-, five-, and six-vertex sextic relaxations defined in Section~\ref{sec:moments}. In the five-vertex relaxation, the eight blocks obtained after structural facial reduction have optimal common spectral margin zero. Removing their exact common kernels yields a strictly feasible point in the reduced blocks. The six-vertex pencil also has a strictly feasible point on its structural face with the selected equality family.
\end{theorem}

Exact arithmetic certifies both three-fixed-line exclusions and the feasible linear functionals on the truncated polynomial spaces.

\section{Gram equations and unitary automorphisms}\label{sec:unitary}

Let $F=[f_1\ \cdots\ f_9]$ be a matrix of unit frame vectors, and use the convention $G_{ij}=f_i^*f_j$ for $G=F^*F$. Set
\[
 \lambda=\frac94,\qquad \mu=\frac{\sqrt{10}}8,\qquad
 \beta=\frac{\lambda-2}{\mu}=\frac{\sqrt{10}}5.
\]
The exact Gram conditions are
\begin{equation}\label{eq:gram}
 G=G^*,\qquad G_{ii}=1,\qquad |G_{ij}|=\mu\quad(i\ne j),\qquad G^2=\lambda G.
\end{equation}
Indeed, $FF^*=\lambda I_4$ follows from tightness and the trace, and applying $\tr(G^2)=\lambda\tr G$ yields $\mu^2=5/32$. Conversely a matrix satisfying \eqref{eq:gram} has eigenvalues in $\{0,\lambda\}$ and trace nine. It is positive semidefinite of rank four and factors as the Gram matrix of the required ETF. The signature matrix $S=(G-I)/\mu$ satisfies
\begin{equation}\label{eq:signature}
 S^2=8I+\beta S.
\end{equation}

A unitary fixing the nine lines is scalar. Its eigenvectors from distinct eigenspaces are orthogonal, whereas $|f_i^*f_j|=\mu>0$ for $i\ne j$. All frame vectors therefore have the same eigenvalue, and they span $\C^4$. Hence the projective unitary automorphism group acts faithfully on the nine lines. A permutation of prime order $p$ has a unitary lift $U$ with $U^p=I$, obtained by rescaling its original lift by a scalar of unit modulus. By Cauchy's theorem it suffices to exclude the following nine cycle types.
\begin{center}
\begin{tabular}{@{}cll@{}}
\toprule
Order & Cycle types & Obstruction\\
\midrule
$2$ & $2\,1^7,\ 2^2 1^5,\ 2^3 1^3,\ 2^4 1$ & Trace and parity\\
$3$ & $3\,1^6,\ 3^2 1^3$ & Fixed lines in a small eigenspace\\
$3$ & $3^3$ & Fourier block correlation\\
$5$ & $5\,1^4$ & Full cyclic span\\
$7$ & $7\,1^2$ & Full cyclic span\\
\bottomrule
\end{tabular}
\end{center}

\paragraph{Order two.}
The lift $U$ is Hermitian and has trace $2r-4$ for $0\le r\le4$. Tightness gives
\[
 \lambda(2r-4)=\sum_j f_j^*Uf_j=a+2\mu b,
\]
where $a$ is the sum of the signs on the fixed lines. A swapped pair contributes two equal terms because $U^2=I$. Each is real since $U=U^*$, and each has modulus $\mu$, so the pair contributes $\pm2\mu$. Thus $b$ is an integer. In all four cycle types $a$ is odd with $|a|\le7$. Irrationality of $\mu$ implies $b=0$. The only integer values of the left side are $-9,0,9$, giving a contradiction in each type.

\paragraph{Orders three, five, and seven with fixed lines.}
For an orbit of length $m$, let $K$ be the Gram matrix of $f,Uf,\ldots,U^{m-1}f$. Cauchy--Schwarz on its nonzero eigenvalues gives
\begin{equation}\label{eq:orbit-rank}
 \rank K\ \ge\ \frac{(\tr K)^2}{\tr(K^2)}
 =\frac{m}{1+(m-1)\mu^2}.
\end{equation}
The projections of the orbit vectors to an eigenspace of $U$ are proportional. In order three the lower bound $16/7$ forces an orbit rank of three. All three cube roots of unity occur, so each eigenspace has dimension at most two. In types $3\,1^6$ and $3^2 1^3$ at least three fixed lines occur. Because they are pairwise nonorthogonal eigenlines, they lie in one eigenspace. Their Gram matrix also has rank at least $16/7$, contradicting its rank at most two. For the five- and seven-cycles, the lower bounds are $40/13$ and $112/31$, respectively. In both cases the orbit spans $\C^4$, so all eigenspaces are one-dimensional. Each type has at least two distinct nonorthogonal fixed lines and these must again belong to one eigenspace. This excludes both types.

\paragraph{Order three without fixed lines.}
Label the lines $(a,x)$ with $a,x\in\{0,1,2\}$, choosing representatives $f_{a,x}=U^x f_{a,0}$. Then $G_{(a,x),(b,y)}=g_{ab}[y-x]$, with indices in brackets taken modulo three. For $\omega=e^{2\pi i/3}$, its Fourier blocks are
\[
 \widehat G_k[a,b]=\sum_{t=0}^2 g_{ab}[t]\omega^{kt},\qquad F_k=\widehat G_k/\lambda.
\]
The transform uses columns $\delta_{ab}\omega^{kx}/\sqrt3$. Thus \eqref{eq:gram} makes $F_k$ orthogonal projections whose ranks add to four. Each block has diagonal entries at least $(1-2\mu)/\lambda>0$ and hence is nonzero. Multiplying each $f_{a,x}$ by $\omega^{mx}$ shifts the frequency labels, so we may assume the ranks are $(2,1,1)$ in that order.

For a triple of matrices $B=(B_0,B_1,B_2)$ define
\[
 \mathcal M(B)=\sum_{k=0}^2 B_{k+1}\odot\overline{B_k}
   -\diag\bigl((B_1+B_2)_{00},(B_1+B_2)_{11},(B_1+B_2)_{22}\bigr),
\]
where $\odot$ is entrywise multiplication and frequency indices are cyclic. Fourier orthogonality gives
\[
 \sum_k F_{k+1}[a,b]\overline{F_k[a,b]}
  =\frac{3}{\lambda^2}\sum_t |g_{ab}[t]|^2\omega^t.
\]
The right side is zero off the diagonal and equals $3(1-\mu^2)/\lambda^2=1/2$ on the diagonal. Also $\sum_k F_k[a,a]=3/\lambda=4/3$. It follows that $\mathcal M(F)$ is diagonal and
\[
 \mathcal M(F)_{aa}=F_0[a,a]-\frac56
 \le \frac{1+2\mu}{\lambda}-\frac56=-\frac{7-2\sqrt{10}}{18}<0.
\]
On the other hand, write $F_0=I-uu^*$, $F_1=vv^*$, and $F_2=ww^*$ for unit vectors. With $p=v\odot\bar u$, $q=w\odot\bar v$, and $r=u\odot\bar w$, expansion gives
\[
 \mathcal M(F)=-pp^*+qq^*-rr^*.
\]
A nonzero vector orthogonal to $p$ and $r$ has nonnegative quadratic form, contradicting the negative diagonal. This excludes $3^3$ and completes the unitary part of Theorem~\ref{thm:asymmetry}.

\begin{proposition}[Quantitative free-$C_3$ obstruction]\label{prop:c3residual}
Let $G$ be Hermitian and invariant under the three simultaneous shifts above. Assume the unit-diagonal and modulus conditions in \eqref{eq:gram}. Then
\[
 \|G^2-\lambda G\|_F\ge
 \frac{\lambda^2(7-2\sqrt{10})}{36(2\sqrt3+\sqrt2)}.
\]
\end{proposition}
\begin{proof}
Write $\varepsilon=\|G^2-\lambda G\|_F$ and $\gamma=(7-2\sqrt{10})/18$. Round the eigenvalues of each Fourier block of $G$ to the nearer of $0$ and $\lambda$ to obtain $H$. For a real eigenvalue $s$, the larger of $|s|$ and $|s-\lambda|$ is at least $\lambda/2$. Hence $\eta=\|G-H\|_F\le2\varepsilon/\lambda$. Suppose $\varepsilon<\lambda^2\gamma/[2(2\sqrt3+\sqrt2)]$. This is smaller than both $\lambda^2/12$ and $\lambda(1-2\mu)/2$. The first threshold gives $|\tr H-9|\le3\eta<\lambda/2$, forcing rank four. The second gives $\eta<1-2\mu$, so each rounded block is nonzero. After shifting the frequencies, $P_k=\widehat H_k/\lambda$ have ranks $(2,1,1)$.

Let $\delta^2=\sum_k\|P_k-F_k\|_F^2=\eta^2/\lambda^2$. Every entry of a projection has modulus at most one, as does every entry of $F_k$ by the bounds $(1+2\mu)/\lambda$ on the diagonal and $3\mu/\lambda$ off the diagonal. Taking differences one factor at a time, the three Hadamard products in $\mathcal M$ contribute at most $2\sqrt3\,\delta$ in Frobenius norm. The diagonal subtraction contributes at most $\sqrt2\,\delta$. Thus $\|\mathcal M(P)-\mathcal M(F)\|_2\le(2\sqrt3+\sqrt2)\delta<\gamma$. The entrywise Fourier calculation gives $\mathcal M(F)\preceq-\gamma I$, while the rank decomposition guarantees a nonnegative direction for $\mathcal M(P)$, a contradiction.
\end{proof}

\paragraph{Root-of-unity phases.}
Equation \eqref{eq:signature} also excludes any Gram matrix that can be rephased to have only roots of unity in the off-diagonal signature entries. In such a gauge
\[
 \beta=\sum_{\ell\notin\{i,j\}} S_{i\ell}S_{\ell j}\overline{S_{ij}}
 \qquad (i\ne j)
\]
is an algebraic integer. Its square $2/5$ contradicts rational algebraic integrality.

\section{Antiunitary automorphisms}\label{sec:anti}
Suppose an antiunitary $A$ preserves the line set. Its square is a scalar unitary by Section~\ref{sec:unitary}, so the permutation is an involution or the identity. In either case it has a fixed line. If $Af=c f$ on this line with $|c|=1$, then $A^2 f=\bar c c f=f$, so $A^2=I$. The real fixed space $\{v:Av=v\}$ is a real form of $\C^4$. The two summands of $(v+Av)/2+i(v-Av)/(2i)$ express $v$ as the sum of a fixed vector and $i$ times a fixed vector, and inner products between fixed vectors are real. In an orthonormal basis of this real form, $A$ is coordinate conjugation. We choose real unit representatives on fixed lines and conjugate unit representatives on swapped pairs.

Five fixed lines would give a real Gram matrix $I+\mu J_0$ of order five and rank at most four, with $J_0$ an integral matrix having off-diagonal entries $\pm1$ and zero diagonal. Then $-1/\mu$ is an algebraic integer, contrary to $(1/\mu)^2=32/5$. Thus there are one or three fixed lines, giving types $2^4 1$ and $2^3 1^3$. When there are three, the product of the three fixed-edge signs is invariant under rephasing by real signs. By choosing positive inner products with the first fixed line, we leave one sign $\sigma\in\{+1,-1\}$ on the remaining edge. On each swapped pair we can independently choose the phase of $z_j$ and the conjugate phase of $\bar z_j$ so that its inner products with the first fixed vector are both $\mu$.

\subsection{Real Gram coordinates}

Replace each pair $(z_j,\bar z_j)$ by $u_j=\sqrt2\Re z_j$ and $v_j=\sqrt2\Im z_j$. This replaces the two columns by their product with the unitary matrix $2^{-1/2}\begin{psmallmatrix}1&-i\\1&i\end{psmallmatrix}$. The resulting real Gram matrix has the form
\[
 R=\begin{pmatrix}K&C\\ C^T&D\end{pmatrix},\qquad R^2=\lambda R.
\]
Here $K$ is the Gram matrix of the fixed vectors, and $D$ is indexed by the real two-column blocks $(u_j,v_j)$. For $j\ne k$, writing $D_{jk}=\begin{psmallmatrix}a&b\\c&d\end{psmallmatrix}$ gives
\begin{align*}
 |\langle z_j,z_k\rangle|^2 &= \bigl((a+d)^2+(b-c)^2\bigr)/4,\\
 |\langle z_j,\bar z_k\rangle|^2 &= \bigl((a-d)^2+(b+c)^2\bigr)/4.
\end{align*}
The difference is $\det D_{jk}$, so $\det D_{jk}=0$. The same calculation on a diagonal block, with $\|z_j\|^2=1$ and $|\langle z_j,\bar z_j\rangle|^2=\mu^2$, gives $\det D_{jj}=1-\mu^2$. The fixed block $K$ is positive definite, since it has size at most three and $1-2\mu>0$. The Schur complement
\begin{equation}\label{eq:schur}
 Z=D-C^T K^{-1} C
\end{equation}
has rank $4-\rank K$ by congruence to $\operatorname{diag}(K,Z)$ and $\rank R=4$.

\subsection{A rational identity for one fixed line}

Let $e_0$ be the real unit representative of the fixed line. Write $u_j=a e_0+x_j$ and $v_j=y_j$, where $a=\sqrt2\mu$ and $x_j,y_j\in e_0^\perp\cong\R^3$. Set $X=[x_1\ \cdots\ x_4]$, $Y=[y_1\ \cdots\ y_4]$ and $W=[x_1\times y_1\ \cdots\ x_4\times y_4]$. The tight frame operator on $e_0^\perp$ yields $T=0$, where
\[
 T=XX^T+YY^T-\lambda I_3.
\]
The wedge inner products $\langle u_j\wedge v_j,u_k\wedge v_k\rangle$ are $\det D_{jk}$. Decomposing the wedges into $\Lambda^2 e_0^\perp$ and $e_0\wedge e_0^\perp$ therefore yields $E=0$, where
\[
 E=W^TW+a^2Y^TY-bI_4,\qquad b=1-\mu^2=27/32.
\]
In fact $YY^T\preceq\lambda I_3$ gives $Y^TY\preceq\lambda I_4$, so $W^TW\succeq(b-a^2\lambda)I_4$. Since $b-a^2\lambda=9/64>0$ and $W$ has three rows, the equations are inconsistent.

For a polynomial certificate, take $h\in\R^4$ with $Wh=0$ and $h^Th=1$, which is possible for every such $W$. Put $\Delta=b-a^2\lambda=9/64$. Then
\begin{align}\label{eq:fixed1identity}
 -1={}&\frac{a^2}{\lambda\Delta}
 \bigl(\|\lambda h-Y^TYh\|^2+\|X^TYh\|^2\bigr)+(h^Th-1)\notag\\
 &+\frac{h^TEh-(Wh)^T(Wh)-(a^2/\lambda)(Yh)^TT(Yh)}{\Delta}.
\end{align}
To check the identity, set $v=Yh$. The two squared norms sum to $\lambda^2h^Th-\lambda v^Tv+v^TTv$, and $h^TEh-(Wh)^T(Wh)=a^2v^Tv-bh^Th$. Substitution cancels the $v$ terms and gives $-1$. At a solution only the positive multiple of the two squared norms remains. The coefficients of the identity are rational since $a$ occurs only through $a^2=5/16$. The archived certificate encodes these eight squares as rank-one Gram blocks. Its verifier reconstructs the constraints and checks the polynomial identity and the block ranks in exact arithmetic (Appendix~\ref{app:replay}).

\subsection{Polynomial certificates for three fixed lines}

Use labels $0,1,2$ for the fixed lines and pair labels $(3,4),(5,6),(7,8)$. For each sign $\sigma$, the following assignments define the entire complex Hermitian matrix $G(x)$, completed by Hermitian transposition. Its diagonal is one and $G_{0j}=\mu$ for $j>0$, with $G_{12}=\sigma\mu$. In the following list, each new phase $z$ or $w$ has the form $\mu(x_{2k}+i x_{2k+1})$ for a new index $k$.
\begin{enumerate}
\item For each left pair index $a\in\{3,5,7\}$ in increasing order, set $G_{a,a+1}=z$.
\item Immediately after the internal edge of that pair, for $i=1,2$ in order set $(G_{i,a},G_{i,a+1})=(z,\bar z)$ using a new phase for each $i$.
\item Still in the loop for $a$, for each $b>a$ among $\{3,5,7\}$ in order set
\[
 (G_{a,b},G_{a+1,b+1},G_{a,b+1},G_{a+1,b})=(z,\bar z,w,\bar w)
\]
using two new phases in the order $z,w$.
\end{enumerate}
There are 15 phases and 30 real coordinates. All solutions satisfy the circle equations $x_{2k}^2+x_{2k+1}^2-1=0$ and therefore lie in $[-1,1]^{30}$. Let $Q$ be the block diagonal matrix with $I_3$ on the fixed columns and the pair transform above on the other columns. Form $R=Q^*G(x)Q$ and $Z$ by \eqref{eq:schur}. The equations $h_i=0$ used in the certificate consist of the circles, every upper-triangular entry of $R^2-\lambda R$, and the minors
\[
 Z_{a,i}Z_{b,j}-Z_{a,j}Z_{b,i},
 \qquad 0\le a<b<6,\quad 0\le i<j<6,\quad (a,b)\le(i,j)
\]
in lexicographic order on the row and column pairs. These are the 120 upper-triangular entries of the matrix of second minors before dropping identically zero equations. A frame must satisfy them by $\rank Z=1$. The lists are independently reconstructed by the verifier. All coefficients belong to $\Q(\sqrt2,\sqrt5)$, with ordered rational basis $(1,\sqrt2,\sqrt5,\sqrt{10})$.

For each sign, the stored polynomial multipliers $m_i$ are rational and the sum-of-squares factors $L_b$ are integer matrices with a common positive denominator $d$. With $v_b$ the archived monomial vector for block $b$, define the residual polynomial
\begin{equation}\label{eq:fixed3residual}
 r=-1-\sum_i m_i h_i -\sum_b \|L_b^T v_b/d\|^2.
\end{equation}
To bound $r$, expand its coefficients exactly in the given field basis. If $N$ is the stored positive interval denominator, enclose each $\sqrt{k}$, $k=2,5,10$, between $\lfloor\sqrt{kN^2}\rfloor/N$ and $(\lfloor\sqrt{kN^2}\rfloor+1)/N$. Rational interval addition with sign-aware multiplication by the four basis coefficients gives an interval for each coefficient of $r$. Summing the larger absolute endpoint over all monomials bounds $|r(x)|$ on the unit box. The rational bounds from the two archived certificates are
\begin{align*}
 B_{+}&=\frac{117823544948222256047655021017435399191}
 {291600000000000000000000000000000000000000000}
 <4.05\times10^{-7},\\
 B_{-}&=\frac{414220049013296592397744354622927415857}
 {777600000000000000000000000000000000000000000}
 <5.33\times10^{-7}.
\end{align*}
At a common zero of the $h_i$, equation \eqref{eq:fixed3residual} gives $r\le-1$, contradicting $|r|\le B_\sigma<1$. The exact argument uses rational field coordinates, interval endpoints, and the integer Gram factors. Appendix~\ref{app:replay} gives the commands to reconstruct these bounds and check the provenance of the equations. These two contradictions and \eqref{eq:fixed1identity} exclude every antiunitary case and complete Theorem~\ref{thm:asymmetry}.


\section{Moment relaxations of the unrestricted Gram equations}\label{sec:moments}
\subsection{Anchoring and invariant moments}

Fix one line as an anchor and rephase the other representatives to make all its incident Gram entries $+\mu$. Index the eight nonanchor vertices by $0,\ldots,7$. Write the signature matrix in \eqref{eq:signature} as
\[
 S=\begin{pmatrix}0&\one^T\\\one&T\end{pmatrix},\qquad
 T_{ii}=0,\quad T_{ij}=x_{ij}+i y_{ij}\quad(i<j),
\]
with $T=T^*$. The 56 real variables satisfy
\begin{equation}\label{eq:anchored}
 x_{ij}^2+y_{ij}^2=1,\qquad T\one=\beta\one,\qquad
 T^2=8I+\beta T-J,\qquad J=\one\one^T.
\end{equation}
The first equations give the diagonal of the last matrix equation. The remaining real equations consist of a circle, real and imaginary row sums, and real and imaginary off-diagonal tightness, together with relabelings of these five types. Conversely, the block equations recover \eqref{eq:signature} and hence \eqref{eq:gram}.

A truncated moment functional $L$ is linear on real polynomials of the given degree limit. The constraints are $L(1)=1$, $L(hm)=0$ for selected equation multiples, and $L(p^2)\ge0$ for all $p$ in a specified monomial span $V$. Here $h$ is an equation polynomial and $m$ is a monomial. The last condition is the positive semidefiniteness of the matrix with entries $L(uv)$ in a basis of $V$. This is the truncated moment construction of Lasserre \cite{lasserre}. The quartic model uses all monomials of degree at most two for $V$ and every equation multiple from \eqref{eq:anchored} of total degree at most four.

Relabeling the eight nonanchor vertices and conjugating all entries preserve \eqref{eq:anchored}. These operations act on monomials by signed permutations, with a sign for each reversal of a $y$-edge and for conjugation of each $y$ factor. We average $L$ over $S_8\times C_2$. A putative frame need only give point evaluations before averaging. Averaging takes place in moment space, irrespective of the frame's projective automorphism group. Invariance permits one coordinate per signed monomial orbit. An orbit whose representative can be mapped to its negative has moment zero. These rules also determine all unused entries of a stored moment matrix from the representative coordinates.

\subsection{Exact reduction of the positivity condition}

For $k=3,4,5,6$, let $V_k$ consist of all monomials of degree at most two and the cubic monomials involving at most $k$ nonanchor vertices. The vertex support of a monomial is the union of the endpoints of its factors, counting each endpoint only once. Thus $V_6$ contains every cubic, while the lower stages include specified subsets. The moment matrices for invariant $L$ commute with the orthogonal signed-permutation action in the monomial basis. This is the setting of symmetry reduction for semidefinite constraints \cite{gatermannparrilo}.

Let $s=\max(4,k)$ for $V_k$ and $s=4$ for the quartic basis. Take $H$ to be the permutation group on the last $8-s$ vertices. Positivity on $V$ is equivalent to positivity on $V^H$. To see this, suppose the full symmetric matrix has a negative eigenvalue and let $P$ be its spectral projection. There is a basis monomial $m$ with $Pm\ne0$. Relabel $m$ to use only the first $s$ vertices. The relabeled monomial is fixed by $H$, and its projection under $P$ remains nonzero and fixed by $H$ because $P$ commutes with the signed group action. This gives a negative direction in $V^H$. The reverse implication follows from restriction.

The transpositions $(0\ 1)$ and $(2\ 3)$ and the conjugation involution commute on $V^H$. Each has two possible eigenvalues, giving eight simultaneous character spaces. Invariance makes the mixed matrix entries between different characters zero. Within a character, a basis consists of signed orbit sums of monomials for the subgroup generated by $H$ and these three involutions. Orbits with inconsistent signs contribute no basis vector. The orbit sums are orthogonal in monomial coefficient space, and their squared norms are the orbit sizes. The moment blocks in these bases are integer linear combinations of the orbit moments. All eight blocks being positive semidefinite is equivalent to the full positivity condition on the chosen $V$.

\subsection{Sextic equations and structural kernels}

The sextic stages retain the quartic equations. They add $L(hm)=0$ for the five types of base equations in \eqref{eq:anchored}, with total degree five or six and multipliers on at most the first $\min(k,4)$ vertices, using representatives for $h$ whose designated row or edge is in this set. Invariance determines their relabelings. The internal summation indices of $h$ range over the full set of eight vertices. We also impose rank equations. All $5\times5$ minors of $32G$ vanish, as do all $4\times4$ minors of $32(G_{\mathrm{tail}}-\mu^2 J)$, where $G_{\mathrm{tail}}$ is the principal Gram submatrix complementary to the anchor. Indeed, the second matrix is a scalar multiple of the Schur complement of $G$ at its anchor, so it has rank three. We take real and imaginary parts, with all multipliers of degree at most one for the Gram minors and at most two for the Schur minors, modulo relabeling. These constraints are necessary even for nonsymmetric frames. The model uses both principal and nonprincipal minors.

For the three- and four-vertex stages we take products $hm\in V_k$, with $h$ ranging over relabelings of the base equations and $m$ over monomials. We project $hm$ into each tail-fixed character subspace and require its moment pairing with every polynomial in that character block to vanish. In the five-vertex stage we add the raw pairings $L((hm)v)=0$ for such products and monomials $v\in V_5$ to the inherited four-vertex equations. Invariance and linearity imply the projected pairings. All such pairings have degree at most six, and the projected directions are finite linear combinations of relabeled equations times monomials. Their point evaluations vanish at each solution of \eqref{eq:anchored}. The independent verifier uses the archived row provenance to check the imposed systems at each stage.

In a block basis $u_j=\sum_{n\in O_j}\epsilon_n n$, the projection of $hm$ has coefficient
\[
 q_j=\frac{\sum_{n\in O_j}\epsilon_n[n](hm)}{|O_j|}
\]
on $u_j$. Thus the linear cross equations $M_b q=0$ place each such $q$ in the kernel of the block matrix $M_b$. This gives a facial reduction \cite{borweinwolkowicz}. Choose a basis of these directions over $\Q(\beta)$ and a pivot set $P_b$ on which the basis has invertible coordinate restriction. The complementary coordinate vectors then span a direct complement. Every quadratic form value on the original block equals its value on the complementary component since the kernel has been imposed. Consequently the reduced block is exactly the principal submatrix on the nonpivot coordinates. In the three- and four-vertex stages we also archive orthogonal complements, computed using the orbit-size weights.

\subsection{Feasible witnesses and matrix sizes}

The exact witness coordinates belong to $\Q(\beta)$. These witnesses satisfy normalization and all the specified linear equations. The original quartic character blocks are checked by exact PSD elimination. We also verify 48 additional quartic identities obtained by changing the anchor where the products on both sides remain within the degree limit. The equality rank including normalization remains 99, and the exact ranks of the eight original quartic blocks are $24,14,10,14,10,14,10,10$. For the three-vertex stage the independent replay checks the original eight blocks and the 763 homogeneous baseline equations. The structural cross equations are checked separately on the same moment vector. For the four- and five-vertex witnesses, exact tests verify the inherited equations and positivity via the block congruences and principal restriction above. Appendix~\ref{app:replay} gives the artifact paths and commands.

\begin{table}[tb]
\centering
\small
\begin{tabular}{@{}lrrl@{}}
\toprule
Relaxation & Cubic monomials & Moment coordinates & Block orders after structural reduction\\
\midrule
Quartic & 0 & 107 & $72,45,34,38,34,38,25,24$\\
Three-vertex sextic & 2{,}576 & 2{,}547 & $74,55,46,50,46,50,35,39$\\
Four-vertex sextic & 14{,}056 & 2{,}716 & $128,99,86,98,86,98,77,75$\\
Five-vertex sextic & 27{,}496 & 2{,}719 & $373,342,284,302,284,302,239,235$\\
Six-vertex sextic & 30{,}856 & 2{,}719 & $886,855,702,721,702,721,579,573$\\
\bottomrule
\end{tabular}
\caption{Sizes of the five relaxations with exact feasible witnesses. The quartic row lists the eight original character blocks. The five-vertex row precedes its second facial reduction.}\label{tab:sizes}
\end{table}

For the four- and five-vertex stages, the change of coordinates $L(m)=\beta^{\deg(m)\bmod2}z_m$ makes the equality system rational after common factors are removed. The affine charts are constructed by rational elimination and checked against all equations. The rank evidence consists of invertible minors and exact identities for the free columns. The four-vertex equality rank including normalization is $2{,}346$, leaving $370$ free coordinates. The linear map from these directions into the eight PSD blocks has rank $97$. Since the remaining directions can be removed without changing the image, the numerical search can operate in 97 dimensions. The five-vertex stage adds 7{,}027 distinct cross equations, has equality rank $2{,}642$ including normalization, and leaves $77$ free coordinates with PSD image rank $75$. Here the extra equations can reduce the image dimension even though the matrix blocks are larger. Exact checks of rank use chart relations and finite-prime minors certified against denominator-cleared rational data. The floating-point barrier iterates are used to locate candidate witnesses. Exact equality and positivity checks establish the assertions in Theorem~\ref{thm:moment}.

\subsection{The common kernel of the five-vertex pencil}

Let $M_b(y)$ be the eight five-vertex blocks on the first reduced face, with $y$ the 77-dimensional affine chart coordinate. In the archive's character order their common kernels have dimensions $(9,6,4,6,4,6,3,3)$. They are verified by multiplication against the constant block and all 77 coefficient blocks. Restricting to complementary principal coordinates gives orders $(364,336,280,296,280,296,236,232)$. On each of these blocks the same exact feasible witness is positive definite. Consequently each common kernel is the full kernel at that witness.

Define the common margin as the supremum of $t$ for which some $y$ satisfies $M_b(y)-tI\succeq0$ for every $b$. Every $y$ has the archived null directions, so $t\le0$. The exact witness attains $t=0$. Positive diagonal rescaling of the blocks preserves this zero optimal margin, since it preserves the existence of common null directions and the PSD witness. After the second facial reduction, strict feasibility on the eight complementary blocks follows from their positive definiteness at the same witness.

\subsection{Six-vertex extension probe}\label{sec:six}

At six vertices the PSD test space contains all $30{,}856$ cubic monomials in the 56 variables. The probe retains the earlier choice of high-degree equality multipliers and adds the six-vertex structural cross equations. Its equality family is the one selected for the probe, with all cubics admitted in the PSD test space. It generates $99{,}821$ cross-product proposals and $7{,}048$ distinct nonzero rows. With the same $2{,}719$ moment coordinates, two additional independent equations increase the equality rank to $2{,}644$ including normalization. The resulting $75$ free directions have PSD image rank $75$. After structural facial reduction the eight block orders are $(886,855,702,721,702,721,579,573)$, with $2{,}105{,}310$ packed rows stored as sparse maps. The estimated peak storage for the blockwise Newton computation is $1{,}283{,}845{,}760$ bytes.

The exact five-vertex witness fails six of the new rows. For example, in the archive's zero-based parity-scaled coordinates, one of the new constraints sets $z_{90}-z_{2617}-z_{2618}$ to zero. At the five-vertex witness this expression equals $-242106776781/12500000000000$. The check fixes all $2{,}719$ inherited coordinates. Thus any six-vertex witness must move at least one coordinate of the archived five-vertex point.

The guarded computation used a 3{,}520 MiB memory limit and a total wall allowance of 1{,}790 seconds. The model assembly, elimination, and numerical search consumed 1{,}186.7 seconds of guarded wall time, with a peak sampled resident set of 768.8 MiB. After 35 Newton steps, the smallest eigenvalue among the eight unshifted, diagonally scaled blocks was approximately $3.06\times 10^{-8}$. The barrier shift at this iterate was about $-0.02095$. The unshifted blocks determine the positivity test for the returned moment vector.

We rounded the 75 free coordinates to rationals with denominator $10^{14}$ and evaluated the exact affine chart. The independent verifier checked the composition with the five-vertex chart, the new equations, and all $2{,}105{,}310$ packed stencils against the reconstructed symmetry bases. Conjugation by the diagonal degree-parity scaling turns each block of the witness into a rational matrix. An integer congruence gives strict diagonal dominance with positive diagonal in each full structurally reduced block. The exact witness therefore proves strict feasibility on the six-vertex structural face. A separate full-rank modular specialization verifies that the eight affine blocks have trivial common kernels. This completes Theorem~\ref{thm:moment}.


\section{Discussion}\label{sec:discussion}
Theorem~\ref{thm:asymmetry} leaves only frames with no nontrivial projective isometries. In particular, a construction that forces a nontrivial cyclic permutation or a conjugation symmetry cannot produce these parameters. The root-of-unity obstruction applies even when a finite-phase ansatz does not force a permutation symmetry.

For a Hermitian matrix with unit diagonal and exact equiangular modulus, the numerical search minimized the tightness residual $\|G^2-\lambda G\|_F$. In the free $C_3$ class, 48 starts per projection algorithm and local refinement gave a best residual of about $0.7168843846$. With 28 independent phases after anchoring, 1{,}000 unrestricted starts and refinement gave residuals in $[0.7168843846433283,0.7168843846433312]$. Each stored endpoint matched the best free-$C_3$ point to a maximum entrywise error below $7\times10^{-8}$ under a search allowing relabeling, rephasing, and complex conjugation.

A Gram matrix satisfying \eqref{eq:gram} would give a construction. A proof of nonexistence by moment infeasibility would require further constraints, for example additional higher-degree equation multiples or PSD constraints. Existence of a complex $\operatorname{ETF}(4,9)$, equivalently of its Naimark complement $\operatorname{ETF}(5,9)$, remains open to our knowledge \cite{fickusmixon}.

\paragraph{Formalization.}
The companion Lean 4 draft contains declarations for the signature change of variables and for the rational identity in \eqref{eq:fixed1identity} in finite real coordinates. The proposed real refutation starts from the lifted equations and a normalized kernel vector. The derivation of those equations from an antiunitary, the remaining symmetry exclusions, and the archived moment verifiers are outside the scope of the draft. Its pinned build instructions are in \pathcode{etf_search/lean/README.md}. The draft awaits a Lean build. The exact replays reported in this article use the independent Python verifiers.


\appendix
\section{Replaying the certificates}\label{app:replay}
\subsection{Antiunitary certificates}

The numerical discovery runs, exact certificates, and independent verifiers are in the companion repository. The three antiunitary certificates and four feasible-moment witnesses below are present at \href{https://github.com/LinkDry/Guess/tree/968d2d5/etf_search}{revision \texttt{968d2d5}}. The commands use PowerShell syntax and run from the root of that checkout. Replay uses Python 3.11 with NumPy, SciPy, SymPy, mpmath, and python-flint. The commands operate directly on archived certificates and moment vectors.

\begin{verbatim}
$s = 'etf_search/src'
$r = 'etf_search/results'
$a = "$r/round3_sdp"
python -B "$s/verify_sdp_cert.py" `
  "$a/identity_fixed1/certificate.json" --fixed1
foreach ($branch in 'fixed3_positive_dense','fixed3_negative_dense') {
  $b = "$a/schur/$branch"
  python -B "$s/verify_sdp_cert.py" "$b/certificate.json" `
    --model "$b/model.json"
}
\end{verbatim}

The first command reconstructs the constraints of \eqref{eq:fixed1identity} and checks all coefficients and eight rank-one Gram blocks. The next two commands independently reconstruct the circle, real Gram tightness, and Schur-minor equations from Section~\ref{sec:anti}. Each checks 179 used equation polynomials against the stored model and recomputes the exact rational residual bound for the corresponding integer factors. Each verifier reports a passed exact proof.

\subsection{Quartic and three-vertex moments}

\begin{verbatim}
$q = "$r/round4_general_4_9/anchored_s8_reanchored"
python -B "$s/verify_sdp_cert.py" `
  "$q/anchored_exact_primal_candidate.json" `
  --model "$q/anchored_model.json"
$p5 = "$r/round5_level3_4_9"
$p6 = "$r/round6_facial_reduction"
$w3 = "$p6/solver_support3_exact_dixon"
$m3 = "$p5/solver_free_support3_scs/solver_model.json"
$f3 = "$p6/solver_support3_scs/solver_face_model.json"
python -B "$s/verify_sdp_cert.py" `
  "$w3/solver_exact_primal_candidate.json" --model "$m3"
python -B "$s/verify_sdp_cert.py" `
  "$w3/solver_exact_primal_candidate.json" --model "$f3"
\end{verbatim}

The quartic verifier checks the 183 homogeneous baseline equations, 48 change-of-anchor identities, normalization, and eight PSD blocks. The first three-vertex check uses the original quartic and sextic equations and original block bases for that stage. The second also checks the 19{,}856 structural cross equations on the same moment vector. The exact PSD rank of each original three-vertex block equals its principal complement order in Table~\ref{tab:sizes}.

\subsection{Four- and five-vertex moments}

\begin{verbatim}
$p7 = "$r/round7_support4_lmi"
$p8 = "$r/round8_support5_lmi"
$f4a = "$p6/solver_support4_clarabel"
$f4 = "$f4a/solver_face_model.json.gz"
$m4 = "$p5/solver_free_support4_minors/solver_model.json"
$f4w = "$p6/solver_support4_clarabel_streamed/solver_face_model.json.gz"
python -B "$s/verify_sdp_cert.py" "$f4" --audit-face --model "$m4"
python -B "$s/verify_sdp_lmi.py" --model "$f4w" `
  --primal-certificate "$p7/barrier_rank97/exact_original_primal.json"
python -B "$s/verify_support5_rank.py" --input "$p8" `
  --prior-face "$f4" --prior-chart "$p7/affine_chart.json.gz" `
  --prior-audit "$f4a/independent_face_verification.json" `
  --common-kernel "$p8/common_kernel.json" `
  --primal "$p8/barrier_rank75/exact_original_primal.json"
\end{verbatim}

The face audit verifies the source of the kernel directions and the coefficient maps for passage to the principal complements. The four-vertex check verifies normalization and all face equations and uses a rational congruence and strict diagonal dominance for the positive definiteness of the reduced blocks. The five-vertex check also regenerates the new cross equations and packed PSD maps, verifies the two ranks of the chart and its image, and checks the common kernel against every affine coefficient block before testing the witness. It uses the constant and homogeneous parts of the affine chart to verify the reductions and lifts the PSD witnesses by the congruence in Section~\ref{sec:moments}. The five-vertex moment feasibility report records the absence of any assumption fixing lower moments.

\subsection{Six-vertex extension probe}

\begin{verbatim}
$p9 = "$r/round9_support6_probe"
python -B "$s/verify_support6_probe.py" --input "$p9" --prior "$p8" `
  --inherited-primal "$p8/barrier_rank75/exact_original_primal.json"
\end{verbatim}

The inherited-point replay reports \texttt{no\_exact\_equality\_extension} and returns status 1 because the exact vector fails the new equality system. It also verifies the new equality rank and the PSD-direction rank. This test fixes the inherited witness and concerns that point alone.


The following replay checks the six-vertex strict witness. It reconstructs the exact moment vector from the archived checkpoint and rebuilds the coefficient stencils before certifying positive definiteness. The saved vector appears in \pathcode{exact_original_primal.json} in the six-vertex archive.
\begin{verbatim}
python -B "$s/verify_support6_probe.py" --input "$p9" --prior "$p8" `
  --solution "$p9/barrier/solver_primal_dual.npz" `
  --audit-stencils --include-certificates
\end{verbatim}

\subsection{Regression tests and Lean draft}

The archived Python regression suite passed all 86 tests with Python 3.11.9, NumPy 2.4.6, SciPy 1.17.1, SymPy 1.14.0, mpmath 1.3.0, python-flint 0.9.0, SCS 3.3.2, Clarabel 0.11.1, and psutil 7.2.2. Run the suite from the repository root in that environment.
\begin{verbatim}
python -B -m unittest discover -s etf_search/tests -v
\end{verbatim}

The companion Lean sources in \pathcode{etf_search/lean} give the signature algebra and the single-fixed-line identity. To attempt a build under the pinned Lean and mathlib versions, run the following in that directory with Elan available.
\begin{verbatim}
lake update
lake exe cache get
lake build ETF49
lake env lean Check.lean
\end{verbatim}
The final command checks the stated examples and prints the declaration axiom dependencies.


\begin{thebibliography}{99}

\bibitem{welch}
L.~R. Welch, Lower bounds on the maximum cross correlation of signals,
\emph{IEEE Trans. Inform. Theory} \textbf{20} (1974), 397--399.
\href{https://doi.org/10.1109/TIT.1974.1055219}{doi:10.1109/TIT.1974.1055219}.

\bibitem{fickusmixon}
M. Fickus and D.~G. Mixon, Tables of the existence of equiangular tight frames,
\href{https://arxiv.org/abs/1504.00253v2}{arXiv:1504.00253v2} (2016).

\bibitem{szollosi}
F. Sz\"oll\H{o}si, All complex equiangular tight frames in dimension 3,
\href{https://arxiv.org/abs/1402.6429v1}{arXiv:1402.6429v1} (2014).

\bibitem{dhillon}
I.~S. Dhillon, R.~W. Heath, Jr., T. Strohmer, and J.~A. Tropp,
Constructing packings in Grassmannian manifolds via alternating projection,
\emph{Exp. Math.} \textbf{17} (2008), 9--35.
\href{https://doi.org/10.1080/10586458.2008.10129018}{doi:10.1080/10586458.2008.10129018}.

\bibitem{goyeneche}
D. Goyeneche and O. Turek, Equiangular tight frames and unistochastic matrices,
\emph{J. Phys. A} \textbf{50} (2017), 245304.
\href{https://doi.org/10.1088/1751-8121/aa6e16}{doi:10.1088/1751-8121/aa6e16}.

\bibitem{gameofsloanes}
J. Jasper, E.~J. King, and D.~G. Mixon, Game of Sloanes: Best known packings in complex projective space,
\href{https://arxiv.org/abs/1907.07848v1}{arXiv:1907.07848v1} (2019).
\href{https://github.com/gnikylime/GameofSloanes#leader-board}{Author-maintained leaderboard}, accessed November 28, 2026.

\bibitem{gap}
M. Fickus, J. Jasper, and D.~G. Mixon,
The Singer--Zauner gap for equiangular tight frames,
\href{https://arxiv.org/abs/2605.28738v1}{arXiv:2605.28738v1} (2026).

\bibitem{chienwaldron}
T.-Y. Chien and S. Waldron,
A characterization of projective unitary equivalence of finite frames and applications,
\emph{SIAM J. Discrete Math.} \textbf{30} (2016), 976--994.
\href{https://doi.org/10.1137/15M1042140}{doi:10.1137/15M1042140}.

\bibitem{lasserre}
J.~B. Lasserre, Global optimization with polynomials and the problem of moments,
\emph{SIAM J. Optim.} \textbf{11} (2001), 796--817.
\href{https://doi.org/10.1137/S1052623400366802}{doi:10.1137/S1052623400366802}.

\bibitem{gatermannparrilo}
K. Gatermann and P.~A. Parrilo, Symmetry groups, semidefinite programs, and sums of squares,
\emph{J. Pure Appl. Algebra} \textbf{192} (2004), 95--128.
\href{https://doi.org/10.1016/j.jpaa.2003.12.011}{doi:10.1016/j.jpaa.2003.12.011}.

\bibitem{borweinwolkowicz}
J.~M. Borwein and H. Wolkowicz, Facial reduction for a cone-convex programming problem,
\emph{J. Austral. Math. Soc. Ser. A} \textbf{30} (1981), 369--380.
\href{https://doi.org/10.1017/S1446788700017250}{doi:10.1017/S1446788700017250}.

\end{thebibliography}

\end{document}
